La cuestión rectora de este trabajo es la conservación de una estructura
generada cuando cambia su realización matemática. La Holografía Modular
Triádica produce, mediante operaciones aritméticas levantadas, estados que
retienen simultáneamente orientación, incidencia, frontera y memoria. Su
estructura discreta del continuo reúne las prolongaciones compatibles de
esos estados. La realización arquimediana expresa coordenadas mediante el
refinamiento de cilindros; la realización conjugada conserva relaciones
excepcionales de incidencia. La genealogía precede, por tanto, a los valores
y a las formas en que se los representa. Ésta es la tesis de partida que el
desarrollo formal siguiente expone con sus operaciones y dependencias.

El problema geométrico propio del artículo comienza cuando una salida
positiva de aquella construcción se convierte en separación entre nodos,
peso de una red, coordenada de Plücker, momento de una medida, coeficiente
de energía o dirección de deformación reticular. Para justificar ese paso
se necesita determinar qué información recupera cada representación, cómo
se compara con las restantes y qué conserva su refinamiento. El interés
de una realización positiva aumenta de forma decisiva cuando admite una
inversa marcada: la geometría deja entonces de ofrecer sólo una evaluación
del estado y pasa a conservar la información que utilizan sus lectores.
La marca comprende el orden de las posiciones, la escala total y los datos
enriquecidos que el lector particular requiere.

La aportación central se organiza alrededor de tres operaciones. La primera
es la reconstrucción: la carga y la memoria recuperan las coordenadas
del registro, y los cocientes de menores recuperan sus separaciones.
La segunda es el refinamiento compatible: una partición se subdivide
conservando sus extremos, sus marcas y un archivo exacto del detalle.
La tercera es el transporte de estructuras: una forma, un operador o un
lector se lleva a otra representación mediante un mapa explícito. Estas
operaciones permiten demostrar conservación de residuos, energía efectiva,
momentos y espectros en los dominios concretos donde cada identidad tiene
sentido. El calificativo positivo adquiere así varias realizaciones
matemáticas precisas, articuladas por la reconstrucción del mismo dato.

La geometría positiva posee antecedentes propios que permiten identificar
con precisión el nivel de esta contribución. Las mediciones de frontera
de redes dirigidas parametrizan celdas del Grassmanniano totalmente no
negativo en el trabajo de Postnikov~\cite{postnikov}; Scott describe la
estructura de álgebra de conglomerados del anillo de coordenadas
grassmanniano~\cite{scott}. Arkani-Hamed y Trnka introducen el amplituedro
en la descripción geométrica de amplitudes de la teoría supersimétrica
\(\mathcal N=4\) de Yang--Mills en el límite planar~\cite{amplituhedron}.
La formulación de geometrías positivas caracteriza después sus formas
canónicas por polos logarítmicos y residuos recursivos de
frontera~\cite{positive}. El presente desarrollo compone esos lenguajes
de realización con datos previamente generados por HMT. Su resultado
diferencial es una colección marcada, recuperable y refinable, junto con
los mapas que enlazan su forma positiva y sus lecturas materiales.

El refinamiento amplía además el alcance dimensional. El número de nodos
crece mediante subdivisiones nonádicas y permite construir matrices
positivas de todo rango y dimensión externa finitos compatibles con ese
número. La prueba se apoya en determinantes de Vandermonde y formas de
momentos de rango completo. En rango uno se obtiene la geometría convexa
completa, con cobertura por símplices y residuos en toda dimensión finita.
Para rangos superiores se construye una colección explícita de imágenes
amplituhedrales y se preservan sus banderas marcadas. La cobertura global
de un amplituedro de rango superior constituye un enunciado de alcance
distinto del teorema de existencia y compatibilidad de esta colección; las
pruebas posteriores conservan esa distinción en su formulación.

La dimensión física se introduce mediante los operadores correspondientes.
El calendario trítico distingue rotación elíptica, apertura hiperbólica y
transporte parabólico de memoria. Sus realizaciones hamiltonianas conservan
la forma simpléctica y permiten calcular la transformación de Legendre y
los términos de conexión originados por una coordenatización variable. En el sector
gauge, la escala areal--volumétrica, la medida común y las formas de
reflexión anteceden a la reconstrucción del Hamiltoniano de Yang--Mills.
El argumento espectral se presenta junto a la coercividad uniforme, el
testigo que alcanza el primer nivel positivo y los entrelazadores que lo
preservan. La relación con las formas canónicas se establece mediante sus
lectores y sus residuos, mientras la brecha pertenece al Hamiltoniano y a
su espacio físico. Esta jerarquía sitúa cada resultado en el lugar donde
su demostración produce efectivamente la conclusión.

La exposición incorpora primero la construcción formal necesaria para
obtener el estado y su registro. Desarrolla después las realizaciones
positivas, la extensión dimensional, la energía y la dualidad reticular;
reúne finalmente las realizaciones hamiltonianas y el desarrollo gauge.
Los programas de comprobación acompañan a identidades determinadas y a
controles negativos explícitos. La demostración simbólica conserva el
alcance general; la ejecución exacta comprueba las especializaciones
declaradas y su correspondencia con las fórmulas impresas.
